\begin{textsample}{POS Dim 3 – ap – Score 20.00 – ap/ap\_inf\_4.txt}  \label{ex:f3_pos_183}
3.2.2.
A \textbf{criança} da Educação \textbf{Infantil} A Educação \textbf{Infantil} surge como um direito da \textbf{criança} e dever do Estado em complementação à ação de sua família.
A primeira \textbf{infância} é uma fase muito importante e deve ser tratada como tal, pois estão envolvidos de maneira significativa os aspectos físico, emocional, espiritual e cognitivo, que serão os alicerces para a sua aprendizagem e interação com o mundo físico e social.
No entanto, isso não pode ser levado em conta visando apenas à possibilidade de a \textbf{criança} ser bem sucedida no futuro, mas, principalmente, buscando proporcionar -lhe espaços onde possa viver sua vida de hoje, de maneira plena, pois é a base para o desenvolvimento do indivíduo como um todo.
É necessário analisar como ocorre esse processo educativo por meio das fases na \textbf{infância}, elencando -se neste contexto quais as relações produzidas para as \textbf{crianças} vivenciarem esta etapa dentro de uma \textbf{instituição} educacional.
I - \textbf{Bebê} ( 0 a 1 ano e 6 \textbf{meses} ) A \textbf{infância} é um período de grande desenvolvimento.
O cérebro passa por várias fases de desenvolvimento e o corpo sofre grandes transformações nos primeiros anos de vida.
Todas as dimensões e situações da vida do \textbf{bebê} – relações familiares e sociais, condições de saúde a nível físico e psicológico, comportamento e interação social, nível socioeconômico da família – contribuem e influenciam a forma como se desenvolve a nível cognitivo, emocional e motor ao longo da vida, “ uma vez que cada estágio é caracterizado pela aparição de estruturas originais, cuja construção o distingue dos estágios anteriores ” ( PIAGET, 1978 ).
Entretanto, para estimular apropriadamente as \textbf{crianças} é preciso conhecer cada estágio de desenvolvimento pelas quais passará, sendo esses estágios, nos estudos de Piaget ( 1978 ), divididos em sensório-motor até 2 anos, pré-operatório dos 2 aos 6/7 anos, operatório concreto dos 6/7 aos 11/12 anos e o estágio das operações formais dos 11/12 até a vida \textbf{adulta}.
Entende -se que, o \textbf{bebê} progride por etapas cronológicas e sequenciais, isto é, vai atingir uma etapa num intervalo de tempo e a mesma acontece com o suporte de uma que aconteceu anteriormente.
Por exemplo, um \textbf{bebê} começa a andar entre os 8 e os 18 \textbf{meses}.
Não o faz enquanto a sua capacidade física não o permitir.
O primeiro ano de vida é fundamental para a saúde cognitiva, pois durante esse período há o maior e mais rápido \textbf{crescimento} do cérebro e, por isso, o mais suscetível é o aumento das conexões neuronais, qual depende da riqueza de experiências \textbf{sensoriais} que o \textbf{bebê} tem.
Os jogos de exercícios para Piaget ( 1978 ), mostra -se como um meio para o \textbf{estímulo} e desenvolvimento deste primeiro estágio, uma vez que o \textbf{bebê} ou \textbf{criança} pode \textbf{brincar} sozinha e não precisa de regras.
Essas experiências moldam a arquitetura do cérebro e os leva a praticar diversas atividades ( mentais ou biológicas ) onde poderá diferenciar os objetos, seu próprio corpo e o ambiente adaptando -se ao meio, Piaget e Inhelder ( 1995 ).
Assim, percebe -se um conjunto de elementos responsáveis pela desenvoltura e aprendizagem do \textbf{bebê}. 0 a 3 \textbf{meses} - Objetos coloridos para ele(a) seguir com os olhos e alternância de \textbf{estímulo} auditivo ; - Estimulação do diálogo ( vínculo mamãe-bebê ) ; - Atividades \textbf{sensoriais} : \textbf{texturas} macias ; - Sorriso social ( estimular o sorriso do \textbf{bebê} ). 4 a 6 \textbf{meses} - O \textbf{bebê} já começa a distinguir expressões sociais, \textbf{sons} como “ papa ”, “ mama ” ; - Movimento dos braços e das \textbf{mãos} com chocalhos ; - Estimulação do sentar do \textbf{bebê} para trabalhar o tônus muscular ; - Na atividade \textbf{sensorial} - inserir novas \textbf{texturas}, \textbf{livros} \textbf{sensoriais}. 7 a 9 \textbf{meses} - \textbf{Brinca} de “ achou ” ( lembra -se do famoso esconde-esconde? ) ; - O \textbf{bebê} já pega objetos, portanto, estimular o engatinhar ; - Na atividade \textbf{sensorial} : trabalho com \textbf{texturas} ásperas e macias, cores, formas, quantidades variadas ; - É importante sempre inserir nas atividades \textbf{texturas} diferentes. 10 a 11 \textbf{meses} - Estimulação do andar ; - Estimulação do falar = \textbf{repetir} palavras, diálogos, desenhos animados, histórias ; - Objetos com \textbf{texturas} variadas. 12 \textbf{meses} - Caminhar com \textbf{ajuda} - explorar novos ambientes ; - Uso da caixa mágica para encontrar objetos ; - A \textbf{criança} \textbf{consegue} falar em média 4 a 5 palavras ( frases curtas ), mediante \textbf{estímulo} do diálogo e histórias ; - Utilização da música para estimular a linguagem ; - Colocar objetos distantes para estimular o andar ; - Atividades de encaixe ; - \textbf{Brinquedos} que estimule a psicomotricidade. 13 a 18 \textbf{meses} - Iniciação da descoberta corporal – atividades para promover a consciência corporal ; - Circuitos psicomotores ( a \textbf{criança} já corre e anda.
Lembra -se da \textbf{brincadeira} “ amarelinha ”? ) ; - Apresentação de regras sociais ; - Trabalho com cores, formas, \textbf{texturas}.

% matched lemmas: adulto, ajuda, bebê, brincadeira, brincar, brinquedo, conseguir, crescimento, criança, estímulo, infantil, infância, instituição, livro, mão, mês, repetir, sensorial, som, textura
\end{textsample}
